\documentclass[11pt]{article}
\usepackage[margin=1in]{geometry}
\usepackage{amsmath,amssymb}
\usepackage{graphicx}
\usepackage{booktabs}
\usepackage{hyperref}
\usepackage{microtype}
\usepackage{parskip}
\usepackage{xcolor}
\usepackage{enumitem}

\setlength{\parskip}{6pt}
\setlength{\parindent}{0pt}

\title{\textbf{Adaptive $\sigma$ for Perturbed Optimizers}\\
\large Advisor Meeting --- March 2026}
\author{Kyle Heuton}
\date{}

\begin{document}
\maketitle
\thispagestyle{empty}

\tableofcontents
\vspace{1em}
\hrule
\vspace{1em}

% ============================================================
\section{Motivation}
% ============================================================

The \textbf{perturbed optimizer} (Berthet et al.\ 2020) approximates the
decision-focused gradient by Monte Carlo: draw $N$ perturbations
$\epsilon_n \sim \mathcal{N}(0, \sigma^2 I)$, solve the CO problem
for each perturbed cost $\hat{y} + \epsilon_n$, and use the resulting
decision spread as a surrogate gradient.
The key hyperparameter is $\sigma$, the perturbation magnitude, which
controls an exploration--exploitation tradeoff:
\begin{itemize}[itemsep=2pt]
  \item $\sigma$ too \emph{small}: nearly all $z_n = z_0$ (the unperturbed
        solution) $\Rightarrow$ zero gradient.
  \item $\sigma$ too \emph{large}: decisions become random
        $\Rightarrow$ high-variance, uninformative gradient.
\end{itemize}

The critical scaling issue is that $\sigma$ must track $\|\hat{y}\|$
(predicted cost magnitude). In practice the dense prediction model's
coefficient norm drifts \textbf{100--1000$\times$} during training,
so a fixed $\sigma$ falls off the useful range after a few dozen epochs.
Polynomial models are more stable ($3$--$4\times$ drift), so this problem
is most acute for dense MLPs.

\textbf{Goal}: design $\sigma$ controllers that automatically maintain
useful gradient signal throughout training, without per-run tuning.

% ============================================================
\section{Exploration Metrics}
% ============================================================

We need a scalar that quantifies ``how much exploration is happening''
so we can close a feedback loop around $\sigma$.

\subsection{Rank Change Rate (RCR) --- binary decisions only}

For binary ($0/1$) solutions with sparsity $K$ (e.g.\ top-$K$ knapsack),
RCR measures the fraction of perturbed solutions that differ from the
unperturbed baseline, normalised by the maximum possible Hamming distance:
\begin{equation}
  \mathrm{RCR} = \frac{1}{N} \sum_{n=1}^{N}
    \frac{\|z_n - z_0\|_1}{2K},
  \label{eq:rcr}
\end{equation}
where $z_0$ is the solution under the unperturbed cost $\hat{y}$ and
$z_n$ is the solution under $\hat{y} + \epsilon_n$.
RCR $= 0$ means no exploration; RCR $= 1$ means complete randomisation.

\textbf{Connection to your comment}: You suggested that among instances
where $y$ and $\hat{y}$ are far apart, $\sigma$ should produce nontrivial
gradients for $\geq 50\%$ of perturbations.  This is precisely RCR
$\geq 0.5$ applied \emph{conditional} on a margin threshold $|y - \hat{y}|
> \delta$.  Our global (unconditional) RCR sweet spot of $0.3$--$0.4$ is
consistent: when only a subset of instances have large residuals at any
training step, global RCR $< 0.5$ while conditional RCR $\approx 0.5$ for
the ``active'' subset.

\textbf{Limitation}: RCR is undefined for non-binary decisions
(portfolio weights, matching edge probabilities) and behaves adversarially
if the solution sparsity $K$ changes across instances.

\subsection{OCV\_Y: Objective Coefficient of Variation}

To generalise beyond binary decisions, we measure the normalised spread
of \emph{true} objective values across the $N$ perturbed solutions:
\begin{equation}
  \mathrm{OCV\_Y} = \frac{\mathrm{std}_n\!\left(Y \cdot z_n\right)}
    {|Y \cdot z_0| + \varepsilon},
  \label{eq:ocvy}
\end{equation}
where $Y$ denotes the \emph{true} (observed) cost coefficients and
$z_n$ are the decisions under the perturbed \emph{predicted} costs.
Key properties:
\begin{itemize}[itemsep=2pt]
  \item \textbf{Monotone in $\sigma$}: larger $\sigma$ produces more
        diverse $z_n$, wider spread in $Y \cdot z_n$.
  \item \textbf{Problem-agnostic}: works for any CO problem;
        uses true objectives, not solution structure.
  \item \textbf{RCR equivalence}: on cubic (binary), the correlation
        between OCV\_Y and RCR is $r = 0.996$, so they are essentially
        interchangeable as control targets for binary problems.
\end{itemize}
The controller holds OCV\_Y $\approx \tau$ (a target) via proportional updates.
Both $z_0$ (unperturbed solution) and $z_n$ (perturbed solutions) are
already computed during the forward pass, so OCV\_Y costs no extra solver calls.

\subsection{FracImproving --- convergence diagnostic}

\begin{equation}
  \mathrm{FracImproving} = \frac{1}{N} \sum_{n=1}^{N}
    \mathbf{1}\!\left[Y \cdot z_n > Y \cdot z_0\right].
  \label{eq:fracimproving}
\end{equation}

FracImproving is the fraction of perturbed solutions that beat the
unperturbed solution under \emph{true} costs.
It is \textbf{not} monotone in $\sigma$ and is not suitable as a control
target. However, it is a useful convergence diagnostic: as $z_0 \to z^*$
(near-optimal), fewer perturbed solutions can improve on it, so
FracImproving $\to 0$. In practice:
\begin{itemize}[itemsep=2pt]
  \item \textbf{Long plateau then cliff}: learning rate too high;
        gradient overshoots.
  \item \textbf{Gradual smooth decay}: appropriate learning rate;
        model is converging properly.
\end{itemize}

% ============================================================
\section{Sigma Controllers}
% ============================================================

\subsection{Global Proportional Controller}

At the end of each epoch, aggregate OCV\_Y from all minibatches in that
epoch (accumulating $\hat{y}$, $z_0$, and $z_n$ buffers), then update:
\begin{equation}
  \sigma_{t+1} = \sigma_t \cdot
    \mathrm{clip}\!\left(\frac{\tau}{\mathrm{OCV\_Y}_t},\;
    \frac{1}{s},\; s\right),
  \label{eq:global-update}
\end{equation}
where $\tau$ is the target OCV\_Y and $s > 1$ is a \texttt{max\_step}
parameter that limits the per-epoch change to a factor of $s$
(prevents runaway).

A single $\sigma$ is shared across all training instances.

\subsection{Per-Instance Proportional Controller}

Maintains a $\sigma$ vector of length $N_\mathrm{train}$
(one value per instance). Each training instance $i$ in a minibatch
has its own OCV\_Y computed from the $N$ perturbed solutions for that
instance. The update applies EMA smoothing:
\begin{equation}
  \sigma_i^{(t+1)} = (1 - \alpha)\,\sigma_i^{(t)}
    + \alpha \cdot \sigma_i^{(t)} \cdot
    \mathrm{clip}\!\left(\frac{\tau}{\mathrm{OCV\_Y}_{i,t}},\;
    \frac{1}{s},\; s\right),
  \label{eq:pi-update}
\end{equation}
where $\alpha$ is an EMA weight (default $0.05$).

\textbf{Motivation}: instances where the K-th decision is tight (small
margin between selected/rejected items) need more perturbation to explore
alternative solutions. Per-instance $\sigma$ allows these ``hard''
instances to use a larger perturbation while ``easy'' instances converge
with a smaller one.

\textbf{Empirical result on cubic}: at epoch 50, small-margin instances
receive $\sim 30\%$ larger $\sigma$ (Spearman $r \approx -0.28$ between
$\sigma_i$ and margin). This pattern is structural, not noise --- it
emerges around epoch 50 and is stable.

% ============================================================
\section{Results}
% ============================================================

\subsection{Cubic: $\sigma$ sweet spot and LR interaction}

The cubic task (top-$K$ selection, polynomial DGP) gives a clean
testbed because the polynomial prediction model is near-optimal.

\begin{figure}[h]
  \centering
  \includegraphics[width=0.85\textwidth]%
    {fig_rcr_sweetspot.png}
  \caption{\textbf{RCR sweet spot on cubic.}
    Fixed-$\sigma$ sweep showing the inverted-U relationship between
    $\sigma$ and validation regret.  RCR $0.3$--$0.4$ coincides with
    the best performance, consistent with the interpretation that roughly
    half the instances are ``actively exploring'' at any time.}
  \label{fig:rcr-sweetspot}
\end{figure}

\begin{figure}[h]
  \centering
  \includegraphics[width=0.90\textwidth]%
    {fig_ocv_lr_comparison.png}
  \caption{\textbf{OCV\_Y controller + LR interaction (cubic, dense model).}
    At lr$= 5\text{e-}2$ (left), FracImproving shows a long plateau
    followed by an abrupt cliff --- a signature of gradient overshooting.
    At lr$= 1\text{e-}2$ (middle), the model breaks through and reaches
    val\_regret $\approx 0.005$.
    At lr$= 5\text{e-}3$ (right), convergence is slower but stable.
    The OCV\_Y controller holds exploration at the target throughout all
    three LRs; the bottleneck is the learning rate, not $\sigma$ control.}
  \label{fig:ocv-lr}
\end{figure}

Key findings on cubic:
\begin{itemize}[itemsep=2pt]
  \item Dense coeff-norm drifts 100--1000$\times$ during training;
        fixed $\sigma$ falls off the useful range by epoch~50.
  \item The OCV\_Y controller automatically compensates:
        $\sigma$ tracks $\|\hat{y}\|$ and stays in the sweet spot.
  \item LR is a critical co-factor: lr$=1\text{e-}2$ with
        OCV\_Y$=0.066$ reaches val\_regret $=0.005$ (dense),
        compared to $\sim0.11$ with lr$=5\text{e-}2$.
  \item Polynomial model is LR-insensitive (already near-optimal at
        val\_regret $\approx 0.008$--$0.023$ for all LRs).
\end{itemize}

\begin{figure}[h]
  \centering
  \includegraphics[width=0.85\textwidth]%
    {fig_pi_lr_comparison.png}
  \caption{\textbf{Per-instance vs.\ global controller (cubic).}
    Polynomial model at lr$=5\text{e-}2$: per-instance best is
    $0.0001$ at target $0.045$, vs.\ global best $0.008$.
    Dense model at lr$=1\text{e-}2$: per-instance best $0.009$ at
    target $0.206$ --- no significant improvement over global ($0.005$),
    consistent with representation capacity being the binding constraint
    for dense on cubic.}
  \label{fig:pi-lr}
\end{figure}

\clearpage
\subsection{Knapsack: benchmark scale}

We evaluated on the standard knapsack benchmark
(20 items, 200 train / 100 test instances, DGP noise\_width$=0.5$).
We replaced the default Gurobi solver in the inner loop with a
vectorised numpy DP solver (0.1~ms/call vs.\ 6~ms for Gurobi),
which allowed $N=100$ perturbation samples per instance per batch
(vs.\ $N=10$ with Gurobi), giving higher-quality gradient estimates.

\begin{table}[h]
  \centering
  \caption{\textbf{Knapsack benchmark results.}
    Test regret (lower is better) for the MSE baseline and four
    perturbed-optimizer variants.  All runs: 300 epochs, dense model,
    lr$=1\text{e-}2$.  MSE uses Gurobi at test time.}
  \label{tab:knapsack}
  \begin{tabular}{lc}
    \toprule
    Controller & Test Regret \\
    \midrule
    MSE (Gurobi, baseline) & \textbf{0.0640} \\
    \midrule
    Global OCV\_Y (DP, $N=100$)        & 0.093 \\
    Per-instance OCV\_Y (DP, $N=100$)  & \textbf{0.084} \\
    Per-instance Hamming const.        & 0.087 \\
    Per-instance Hamming decay         & 0.096 \\
    \bottomrule
  \end{tabular}
\end{table}

\begin{figure}[h]
  \centering
  \includegraphics[width=0.85\textwidth]%
    {fig_kn_bench_results.png}
  \caption{\textbf{Knapsack benchmark results.}
    Val regret and FracImproving traces across training epochs for the
    four perturbed controller variants, plus the MSE baseline (dashed).
    FracImproving plateaus at $\approx 0.13$ for all runs --- an
    irreducible floor set by the DGP noise (noise\_width$=0.5$);
    this means roughly 13\% of perturbations will always find solutions
    better than $z_0$ under true costs, even at optimality.}
  \label{fig:kn-results}
\end{figure}

Key findings on knapsack:
\begin{itemize}[itemsep=2pt]
  \item Per-instance $\sigma$ outperforms global ($0.084$ vs.\ $0.093$,
        $\sim 10\%$ improvement).
  \item Hamming decay --- annealing $\sigma \to 0$ over training ---
        is \emph{worse} than constant $\sigma$. Reducing exploration late
        in training kills the gradient signal before convergence.
  \item The MSE baseline still outperforms all perturbed variants by
        $\sim 31\%$ ($0.064$ vs.\ $0.084$). This gap is investigated
        in the next section.
\end{itemize}

\subsection{Local minimum diagnosis (Waves 6--7)}

The 31\% gap prompted a systematic investigation:
is the perturbed optimizer finding a different (worse) local minimum,
or is there a structural barrier?

\begin{figure}[h]
  \centering
  \includegraphics[width=0.85\textwidth]%
    {fig_kn_loss_landscape.png}
  \caption{\textbf{Loss landscape interpolation (Wave~6).}
    Linear interpolation between the random-init solution weights
    $\theta_\mathrm{rand}$ and the MSE solution weights $\theta_\mathrm{MSE}$,
    evaluated at 200 grid points.  Both decision regret and perturbed
    objective decrease monotonically along this path --- no barrier exists
    between the two solutions in this 1D slice.}
  \label{fig:landscape}
\end{figure}

\begin{figure}[h]
  \centering
  \includegraphics[width=0.85\textwidth]%
    {fig_kn_warmstart.png}
  \caption{\textbf{Warm-start from MSE checkpoint (Wave~7).}
    Cold init (random weights): test regret $= 0.085$.
    Warm init (MSE weights, then fine-tune with perturbed loss):
    test regret $= \mathbf{0.064}$, matching MSE exactly.
    The perturbed method initialized from MSE does not improve over MSE ---
    it converges to the same solution.  Random init falls into a worse
    local minimum at $\approx 0.085$.}
  \label{fig:warmstart}
\end{figure}

\textbf{Interpretation}: The warm-start result (Figure~\ref{fig:warmstart})
confirms that the 31\% gap is an \emph{optimization} problem, not a
structural limitation of the perturbed loss:
\begin{enumerate}[itemsep=2pt]
  \item The MSE solution is also a local minimum of the perturbed
        objective; gradient-based fine-tuning does not move away from it.
  \item Random initialization falls into a different (worse) local
        minimum at $0.085$.
  \item No barrier exists along the 1D interpolation path
        (Figure~\ref{fig:landscape}), but the full high-dimensional
        landscape can still have local minima --- the interpolation
        rules out a barrier along one path, not globally.
\end{enumerate}

\textbf{Practical implication}: The benchmark already runs MSE pretraining
(\texttt{--n\_ptr\_epochs}) before PnO fine-tuning.  If the fine-tuning
learning rate is too high, the model may escape the MSE basin and fall
into the nearby $0.085$ attractor.  Lowering the fine-tuning LR may
recover the MSE performance while adding the potential for improvement
via decision-focused gradients.

% ============================================================
\section{Open Questions and Next Steps}
% ============================================================

\begin{enumerate}[itemsep=6pt]

  \item \textbf{Conditional RCR (your suggestion)}.
    Compute RCR conditioned on instances where $\|y - \hat{y}\|_2 > \delta$.
    This isolates ``hard'' instances where the prediction is currently wrong
    and gradient signal matters most.  A natural question: does
    optimising for conditional RCR $\approx 0.5$ outperform global OCV\_Y?
    This requires defining $\delta$ (e.g., top 25\% of residuals),
    but could give a sharper, more interpretable target.

  \item \textbf{Fine-tuning LR on knapsack}.
    Does lowering the fine-tuning LR (from $5\text{e-}2$ to $1\text{e-}2$
    or $1\text{e-}3$) after MSE pretraining close the 31\% gap?
    Cubic suggests yes: LR$=1\text{e-}2$ broke the plateau for dense.
    On knapsack, the experiment is: MSE pretrain for $P$ epochs
    $\rightarrow$ fine-tune with perturbed loss at varied LRs.
    If this closes the gap, the conclusion is that the benchmark's
    default LR schedule is the bottleneck, not the loss function.

  \item \textbf{Generalisation to other problems}.
    We have tested cubic and knapsack.  Does the per-instance OCV\_Y
    controller consistently outperform global, and does it outperform MSE,
    on portfolio (continuous), bipartite matching (combinatorial),
    or TSP (permutation)?

  \item \textbf{Theoretical connection: OCV\_Y $\approx$ RCR}.
    On cubic (binary decisions), OCV\_Y and RCR have near-perfect
    correlation ($r = 0.996$).  Is this provable?
    For a top-$K$ selector, OCV\_Y $\propto K^{-1/2}$ times the expected
    Hamming distance between $z_n$ and $z_0$, suggesting a closed-form
    relationship under mild distributional assumptions.

  \item \textbf{What does the perturbed method actually add?}
    Warm-start results show the method doesn't improve over MSE when
    initialised from MSE weights.  Is there any regime (task, model,
    noise level) where the perturbed gradient provides a genuine
    improvement over MSE fine-tuning?  If not, the practical
    recommendation is: run MSE with good LR, skip the perturbed loss.

\end{enumerate}

% ============================================================
\appendix
\section{Implementation Notes}
% ============================================================

\begin{itemize}[itemsep=4pt]

  \item \textbf{RCR equivalence to entropy}: empirically, RCR and
    decision entropy have $r = 0.996$ on cubic.  Both are valid control
    targets for binary decisions; OCV\_Y supersedes both for general CO.

  \item \textbf{DP solver for knapsack}: custom vectorised numpy
    0-1 DP (\texttt{openpto/method/Solvers/heuristic/dp.py}),
    $0.1$~ms/call.  Allows $N = 100$ samples where Gurobi only allows
    $N = 10$ in the same wall-clock budget.  The extra samples
    meaningfully reduce gradient variance.

  \item \textbf{FracImproving floor on knapsack}: DGP noise\_width$=0.5$
    means even the true $z^*$ is beaten by $\approx 13\%$ of perturbations,
    because noise shifts some costs enough to change the optimal solution.
    This sets an irreducible FracImproving floor --- it does not indicate
    a suboptimal solution.

  \item \textbf{Per-instance $\sigma$ vs.\ margin}: the negative
    Spearman correlation ($r \approx -0.28$) between $\sigma_i$ and
    the K-th-item margin emerges at epoch 50.  The controller assigns
    larger $\sigma$ to tight-margin instances automatically, without
    any explicit margin signal.

  \item \textbf{Cosine similarity as diagnostic}: the cosine similarity
    between the perturbed gradient and the finite-difference gradient is
    $\approx 0$ throughout training on cubic.  The method does not
    work by aligning with the FD gradient; it works by smoothing
    the loss landscape.  Cosine similarity is uninformative as a metric.

\end{itemize}

\end{document}
